% Chapter 2: Fermat's Last Theorem and the Modularity Theorem
\let\LocalMacro\relax

\chapter{Fermat's Last Theorem and the Modularity Theorem}

\section{The Problem}

\subsection{Informal}
In 1637, a French mathematician named Pierre Fermat was reading an ancient Greek math book when he scribbled a note in the margin. He claimed that no whole number raised to a power higher than two can be split into two other whole numbers raised to that same power. For example, there is no way to write two perfect cubes that add up to another perfect cube, and the same is true for fourth powers, fifth powers, and so on forever. Fermat said he had a brilliant proof but never wrote it down. Mathematicians spent over 350 years chasing this ghost.

\subsection{Simplified Version}
Individual exponents were settled long ago: Fermat himself proved $n=4$ using an elegant infinite-descent argument. Euler settled $n=3$ in 1770, and Dirichlet and Legendre handled $n=5$ in 1825. The real challenge was proving the statement for \emph{all} exponents at once---which required an entirely new theory connecting elliptic curves to modular forms.

\subsection{Formal}
The equation
\begin{equation}
x^n + y^n = z^n
\end{equation}
has no nonzero integer solutions for $n > 2$.

More deeply, the Taniyama--Shimura--Weil conjecture (now the Modularity Theorem) asserts:

\begin{conjecture}[Modularity Conjecture]
Every elliptic curve over $\mathbb{Q}$ is modular, meaning its $L$-function arises from a modular form.
\end{conjecture}

If this conjecture holds, then Fermat's Last Theorem follows, because a hypothetical counterexample $a^n + b^n = c^n$ would yield an elliptic curve (the Frey curve $y^2 = x(x - a^n)(x + b^n)$) that could not be modular, contradicting the conjecture.

\begin{historicalbox}
Fermat's marginal note became one of the most famous unsolved problems in mathematics. Generations of mathematicians tried and failed. Many claimed proofs, all were wrong. The problem inspired the development of algebraic number theory.
\end{historicalbox}

\subsection{Why It Matters}
Fermat's Last Theorem is not just a curiosity about whole numbers. Proving it required building an entirely new bridge between two distant areas of mathematics---elliptic curves and modular forms---that has since become one of the most powerful tools in number theory. The techniques developed along the way have been applied to problems in cryptography, coding theory, and beyond.

\subsection{What Was Known}
For over 350 years, mathematicians chipped away at special cases. By the 1990s, the theorem had been verified by computer for all exponents up to about four million. Various partial results were proved using algebraic number theory and cyclotomic fields, but a general proof remained elusive. The key insight came in the 1980s when Gerhard Frey and Ken Ribet showed that proving the Taniyama--Shimura conjecture would imply Fermat's Last Theorem.

\section{Mathematical Theory}


\subsection{Prerequisites}
This chapter assumes familiarity with:
\begin{itemize}
\item Elementary number theory (congruences, factorization)
\item Basic algebraic geometry (elliptic curves)
\item Complex analysis (modular forms basics)
\end{itemize}

\subsection{Elliptic Curves}

An elliptic curve is a smooth projective curve of genus 1 with a rational point. In Weierstrass form:

\begin{equation}
y^2 = x^3 + ax + b
\end{equation}

where $4a^3 + 27b^2 \neq 0$. Elliptic curves form an abelian group under a geometric addition law.

\subsection{Modular Forms}

A modular form is a holomorphic function on the upper half-plane satisfying a transformation law under $SL_2(\mathbb{Z})$:

\begin{equation}
f\left(\frac{az+b}{cz+d}\right) = (cz+d)^k f(z)
\end{equation}

for integers $a,b,c,d$ with $ad-bc = 1$. The Fourier expansion of a modular form encodes deep arithmetic information.

\subsection{The Taniyama--Shimura--Weil Conjecture}

\begin{conjecture}[Modularity Conjecture]
Every elliptic curve over $\mathbb{Q}$ is modular, meaning its L-function arises from a modular form.
\end{conjecture}

This conjecture connected two seemingly unrelated areas: elliptic curves (geometry) and modular forms (analysis).

\subsection{The Frey Curve and Ribet's Theorem}

In 1984, Gerhard Frey observed that if Fermat's Last Theorem were false, one could construct an elliptic curve (now called the Frey curve) with bizarre properties \cite{frey1984}:

\begin{equation}
y^2 = x(x - a^n)(x + b^n)
\end{equation}

where $a^n + b^n = c^n$. This curve would be so unusual that it could not possibly be modular.

\begin{theorem}[Ribet, 1986 \cite{ribet1990}]
If the Modularity Conjecture is true, then Fermat's Last Theorem is true.
\end{theorem}

This was the crucial link: proving modularity for all elliptic curves would settle Fermat's Last Theorem.

\section{The Solution}

\subsection{Wiles's Proof (1993--1995)}

Andrew Wiles, a British mathematician at Princeton, announced a proof of the Modularity Conjecture for semistable elliptic curves in 1993 \cite{wiles1995}. This immediately implied Fermat's Last Theorem.

The proof used:

\begin{enumerate}
\item \textbf{Galois representations}: Associating Galois representations to elliptic curves and modular forms.
\item \textbf{Deformation theory}: Studying how Galois representations vary in families.
\item \textbf{Iwasawa theory}: Deep results about the arithmetic of elliptic curves over $\mathbb{Z}_p$-extensions.
\item \textbf{Kolyvagin--Flach Euler systems}: New tools for bounding Selmer groups.
\end{enumerate}

\begin{highlightbox}
Wiles worked in secret for seven years on this proof. He told no one except his wife. On June 21, 1993, he gave a three-hour lecture at the Isaac Newton Institute in Cambridge, England, revealing the proof. The audience sat in stunned silence when he finished.
\end{highlightbox}

\subsection{The Flaw and the Fix}

After peer review, a gap was found in the proof. Wiles and his former student Richard Taylor spent nearly a year trying to fix it. In September 1994, they realized that combining two different approaches to the deformation theory would work. The corrected proof was published in 1995 in two papers:

\begin{itemize}
\item Wiles (1995): ``Modular elliptic curves and Fermat's Last Theorem'' \cite{wiles1995}
\item Taylor \& Wiles (1995): ``Ring-theoretic properties of certain Hecke algebras'' \cite{taylorwiles1995}
\end{itemize}

\subsection{The Complete Proof}

Breuil, Conrad, Diamond, and Taylor (2001) extended Wiles's methods to prove the full Modularity Conjecture for all elliptic curves over $\mathbb{Q}$ \cite{bcdt2001}, completing the proof of Fermat's Last Theorem.

\section{Impact}

\begin{itemize}
\item \textbf{Number theory}: Wiles's methods revolutionized the study of Galois representations and deformation theory.
\item \textbf{Langlands program}: The proof confirmed deep connections between different areas of mathematics predicted by the Langlands program.
\item \textbf{Mathematical culture}: Wiles's story captured the public imagination and demonstrated the power of sustained dedication to a single problem.
\end{itemize}

\section{Exercises}

\begin{exercise}[Basic]
Prove that $x^4 + y^4 = z^4$ has no nonzero integer solutions using Fermat's method of infinite descent. \emph{Hint: If $x^4 + y^4 = z^4$, then $(x^2)^2 + (y^2)^2 = (z^2)^2$, giving a Pythagorean triple.}
\end{exercise}

\begin{exercise}[Intermediate]
Show that the elliptic curve $y^2 = x^3 + ax + b$ is nonsingular if and only if $4a^3 + 27b^2 \neq 0$. \emph{Hint: Compute the partial derivatives and find where they simultaneously vanish.}
\end{exercise}

\begin{exercise}[Challenge]
Explain why Ribet's theorem implies that proving modularity for semistable elliptic curves implies Fermat's Last Theorem. \emph{Hint: A Frey curve from a Fermat counterexample would be semistable.}
\end{exercise}


\begin{exercise}[Computational]
Write a Python/SageMath script to explore a concept from this chapter. See the appendix for starter code.
\end{exercise}

\section{Timeline}

\begin{tabularx}{\linewidth}{lX}
\toprule
\textbf{Year} & \textbf{Event} \\ \midrule
1637 & Fermat writes his marginal note \\ 
1753 & Euler proves $n = 3$ \cite{euler1783} \\ 
1825 & Dirichlet and Legendre prove $n = 5$ \\ 
1839 & Lam\'e proves $n = 7$ \\ 
1984 & Frey observes the connection to elliptic curves \cite{frey1984} \\ 
1986 & Ribet proves the connection to modularity \cite{ribet1990} \\ 
1993 & Wiles announces the proof \cite{wiles1995} \\ 
1994 & Wiles and Taylor fix the gap \cite{taylorwiles1995} \\ 
1995 & The complete proof is published \cite{wiles1995,taylorwiles1995} \\ 
2001 & Breuil--Conrad--Diamond--Taylor prove full modularity \cite{bcdt2001} \\ \bottomrule
\end{tabularx}

\section{Citations}

\begin{enumerate}
\item Wiles, A. (1995). ``Modular elliptic curves and Fermat's Last Theorem.'' Annals of Mathematics, 141(3), 443--551.
\item Taylor, R., \& Wiles, A. (1995). ``Ring-theoretic properties of certain Hecke algebras.'' Annals of Mathematics, 141(3), 553--572.
\item Breuil, C., Conrad, B., Diamond, F., \& Taylor, R. (2001). ``On the modularity of elliptic curves over $\mathbb{Q}$.'' Inventiones Mathematicae, 145(3), 493--541.
\item Ribet, K. (1990). ``On modular representations of $Gal(\overline{\mathbb{Q}}/\mathbb{Q})$ arising from modular forms.'' Inventiones Mathematicae, 100(1), 43--103.
\item Frey, G. (1984). ``Links between stable elliptic curves and certain Diophantine equations.'' Annales Universitatis Saraviensis.
\item Euler, L. (1783). ``De resolutione formulae arithmeticae.'' Novi Commentarii Academiae Scientiarum Petropolitanae.
\end{enumerate}